% ----------------------------------------------------------------------------
%                                  METODOLOGIA
% ----------------------------------------------------------------------------

\chapter{METODOLOGIA}
\label{chap:metodologia}

Este capítulo descreve os procedimentos metodológicos e técnicos adotados para a concepção, o desenvolvimento e a validação do sistema proposto. O objetivo é detalhar o percurso percorrido desde a identificação do problema até a estruturação da solução tecnológica voltada para o manejo comportamental e a comunicação escolar de alunos com TDAH e TOD. Para isso, são apresentados a classificação da pesquisa, o modelo de engenharia de software escolhido, as técnicas de levantamento de requisitos e as tecnologias utilizadas na construção da plataforma.

\vspace{0.5cm}
\section{Descrição Geral da Pesquisa}

A pesquisa delineada neste projeto classifica-se, quanto à sua natureza, como aplicada, visto que busca gerar conhecimentos para aplicação prática na resolução de problemas de comunicação entre família e escola no cenário da neurodivergência. Segundo \citeonline{prodanov2013}, a pesquisa aplicada caracteriza-se pelo interesse na aplicação e nas consequências práticas do conhecimento.

Em relação aos objetivos, trata-se de um estudo de caráter exploratório que, conforme define \citeonline{gil2002}, visa desenvolver e esclarecer conceitos para proporcionar uma visão geral sobre um determinado fato. Neste trabalho, a etapa exploratória consolida-se através da revisão bibliográfica e da análise de sistemas similares, fundamentais para mapear as lacunas tecnológicas do mercado educacional.

Por fim, a abordagem metodológica é qualitativa. O foco central não reside na quantificação de variáveis numéricas, mas na compreensão detalhada da dinâmica escolar e familiar, permitindo a correta tradução das necessidades pedagógicas e de manejo comportamental em requisitos de software.

Cabe ressaltar que, devido às restrições temporais inerentes ao desenvolvimento deste trabalho, o escopo desta publicação delimitou-se à concepção, especificação técnica e prototipagem de alta fidelidade da aplicação web responsiva, não englobando as fases de implementação do código-fonte e testes empíricos com usuários. Contudo, com o intuito de garantir a continuidade e a futura validação prática da solução proposta, foi firmada uma articulação institucional para que o sistema seja submetido a um estudo de caso em ambiente real no Instituto Federal de Brasília (IFB). Essa etapa futura viabilizará a análise do impacto gerado pelos cards de manejo proativo diretamente na rotina dos docentes e na comunicação com os familiares das crianças neurodivergentes.

\vspace{0.5cm}
\section{3.2 Instrumentos de Pesquisa}

    Para o levantamento de requisitos e a compreensão aprofundada do contexto escolar e
familiar de alunos com 
TDAH e TOD, o presente estudo apoia-se em instrumentos metodológicos diversificados. A escolha dessas 
ferramentas fundamenta-se na necessidade de captar as percepções dos usuários, observar a dinâmica do ambiente 
educacional e validar a solução tecnológica proposta. Os instrumentos definidos para esta pesquisa são:
\vspace{0.3cm}

No contexto deste trabalho, as entrevistas semiestruturadas, fundamentadas em \citeonline{gil2002}, permitiram explorar de maneira flexível as percepções dos participantes, combinando perguntas abertas com tópicos previamente definidos. Isso favoreceu respostas detalhadas e espontâneas, sendo essenciais para captar as necessidades comunicativas e as expectativas de familiares e educadores em relação ao manejo comportamental. Além disso, a observação participante e as visitas em campo viabilizaram a análise do comportamento dos sujeitos em seu ambiente natural, registrando interações e padrões cotidianos \citeonline{creswell2014}. Esse instrumento foi fundamental para entender a dinâmica real da sala de aula, as estratégias já utilizadas pelos professores e as reações espontâneas no suporte às neurodivergências.

Paralelamente a isso, a aplicação de um questionário online, estruturado com questões abertas e fechadas via Google Forms, atuou como um recurso decisivo para mapear as necessidades técnicas e as expectativas dos usuários. A coleta de dados através desse questionário permitiu identificar não apenas as demandas gerais, mas também as principais "dores" e limitações enfrentadas pelos docentes no seu cotidiano. A análise das respostas revelou dois pontos cruciais que redefiniram o escopo do projeto.

Em primeiro lugar, os educadores manifestaram a preocupação de que a plataforma pudesse tornar-se apenas mais um meio de registro passivo, funcionando como uma versão digital do tradicional "livro de atas" ou diário de turma, onde os acontecimentos são documentados, mas sem a oferta de soluções práticas. Para mitigar esta questão, definiu-se que o sistema deveria assumir uma postura proativa, fornecendo sugestões, exemplos de manejo comportamental e estratégias de intervenção baseadas em evidências para auxiliar no suporte a alunos com TDAH e TOD.

Em segundo lugar, o levantamento expôs um obstáculo infraestrutural: a recusa de alguns docentes em utilizar os seus celulares e redes de dados pessoais para o trabalho escolar. No entanto, constatou-se que as escolas estão equipadas com computadores para uso institucional. Para contornar essa limitação técnica e garantir a adesão ao sistema, decidiu-se que o projeto seria desenvolvido sob a arquitetura de uma Aplicação Web Responsiva. Esta abordagem permite que a plataforma seja acessada diretamente através de um navegador \textit{(browser)} em qualquer computador da escola, sem a necessidade de instalar aplicativos. Em simultâneo, a responsividade garante que a interface se adapte de forma automática e otimizada às telas dos celulares dos pais, permitindo-lhes acompanhar as ocorrências com total mobilidade.

Portanto, diante dessas definições, avançou-se para a prototipagem de alta fidelidade. A interface desenvolvida (construída no Figma) atuará como um instrumento visual de validação, permitindo que os usuários visualizem as funcionalidades propostas e forneçam \textit{feedbacks} iniciais sobre a usabilidade, a clareza visual e a acessibilidade cognitiva do sistema.

\vspace{0.5cm}
\section{Coleta e Análise de Dados}

Nesta etapa, define-se como os dados coletados pelos instrumentos de pesquisa serão processados e interpretados, a partir do uso experimental e da validação dos protótipos do sistema nas dependências do Instituto Federal de Brasília. Durante as sessões de apresentação da ferramenta, serão registradas as reações espontâneas dos educadores e familiares, possíveis dificuldades de navegação na interface web, preferências de layout e o nível de compreensão das funcionalidades propostas. As entrevistas semiestruturadas complementarão esses dados, fornecendo impressões detalhadas dos professores e responsáveis sobre a utilidade da plataforma no cotidiano.

As sessões de observação e as visitas em campo permitirão acompanhar a dinâmica escolar em tempo real, registrando os comportamentos, as estratégias de interação adotadas pelos docentes no manejo do TDAH e do TOD, e as possíveis dificuldades enfrentadas no registro tradicional de ocorrências.

O volume de dados qualitativos obtidos (transcrições de entrevistas e anotações de campo) será organizado e interpretado por meio da Análise de Conteúdo \citeonline{bardin2011}. Esse método permitirá a identificação de categorias temáticas, padrões de comportamento, gargalos na comunicação escola-família e aspectos cruciais de usabilidade. Essa abordagem estruturada possibilitará avaliar a eficácia conceitual do aplicativo web, fornecendo subsídios e métricas concretas para o refinamento contínuo da ferramenta e a consolidação dos requisitos do software.